	% 封面示例————“双学位学士”以外的学位论文
%	\uestccover[<学院名称排版风格>]{<论文题目>}{<学科专业>}{<学号>}{<作者姓名>}{<指导教师>}{<教师职称>}{<学院>} % 单学位学士、研究生通用
% \uestccover{}{}{}{}{}{}{}  % 空封面

\uestccover{城市时空计算中的大规模轨迹数据高效处理研究}
            {计算机科学与技术}
            {202211081351}
            {李\hspace{1em}科}
            {商  烁}
            {教授}%{终南山古墓派}
            {计算机科学与工程学院(网络空间安全学院)}  % 默认将过长的学院名称压缩至下划线宽度

% 中文扉页， 仅研究生用
\DegLv{}  % 清空文档类选项设置的<申请学位级别>
% \uestczhtitlepage  % 空白中文扉页

\ClsNum{TP399}  % \ClsNum{<分类号>}
\ClsLv{公开}  % \ClsLv{<密级>}
\UDC{004.9}  % \UDC{<UDC号>}
\DissertationTitle{城市时空计算中的大规模轨迹数据高效处理研究}  % \DissertationTitle{<题名>}
\Author{李科}  % \Author{<作者姓名>}
\Supervisor{商烁}{教授}{电子科技大学}{成都}  % \Supervisor{<指导教师>}{<职称>}{<单位名称>}{<单位地址>}
% 副导师信息， 无则注释
%  \AssociateSupervisor{陈力思}{教授}{电子科技大学}{成都}
\DegLv{博士}  % \DegLv{<申请学位级别>}， 该信息由文档类选项自动确定， 需修改默认内容时使用， 否则注释即可
\Major{计算机科学与技术}  % \Major{<学科专业>}
\Profield{}  % \Profield{专业学位领域代码}， 此为专业学位独有， 学术学位用户注释即可
\Date{2026年3月18日}{2026年5月26日}  % \Date{<论文提交日期>}{<论文答辩日期>}
\Grant{电子科技大学}{2026年6月}  % \Grant{<学位授予单位>}{<学位授予日期>}
\Reviewer{姜仁河 \hspace{1em} 教授}{}  % \Reviewer{<答辩委员为主席>}{<评阅人>}
\uestczhtitlepage
% \uestczhtitlepage[compress]  % 在compress选项下，“中文扉页”将自动压缩过长职称的字宽， 使之不超出官方模板为该内容设定的长度


% 英文扉页， 仅研究生用
% \uestcentitlepage{<文题>}{<专业>}{<学号>}{<作者>}{<导师>}{<副导师>}{<学院>}， 若无副导师， 则将“<副导师>”参数留空即可
% \uestcentitlepage{}{}{}{}{}{}{}  % 空白英文扉页
\uestcentitlepage{Efficient Processing of Large-Scale Trajectory Data in Urban Spatiotemporal Computing}
                {Computer Science and Technology}
                {202211081351}
                {Li Ke}
                {Prof.Shang Shuo}{}
                {School of Computer Science and Engineering (School of Cyber Security)}

% 独创性声明：[<签名宽度>]{<日期>}{<作者签名图片1>}[<作者签名图片2 默认值为作者签名图片1>]{<导师签名图片>}， 仅研究生用
% \declaration{}{}{}  % 空白独创性声明

\declaration[2cm]{2026年5月26日}{authsign}{spvrsign}

% \declaration[2cm]{}{}[]{spvrsign}